\documentclass[11pt,a4paper]{article}

\usepackage[utf8]{inputenc}
\usepackage[T1]{fontenc}
\usepackage[spanish,es-nodecimaldot]{babel}
\usepackage{lmodern}
\usepackage{geometry}
\usepackage{enumitem}
\usepackage{hyperref}
\usepackage{graphicx}

\geometry{margin=2.4cm}
\setlength{\parindent}{0pt}
\setlength{\parskip}{0.45em}
\setlist[itemize]{topsep=2pt,itemsep=1pt,leftmargin=1.4em}

\hypersetup{
	colorlinks=true,
	linkcolor=black,
	urlcolor=blue,
	pdftitle={Oferta de TFG: Interfaz gráfica para una máquina universal de ensayos}
}

\title{Generación de Estereogramas Lunares Mediante Paralaje por Rotación Terrestre}

\begin{document}

\maketitle

\textbf{Titulación:} Grado en Ingeniería Informática\\
\textbf{Centro:} Universidad de Almería\\
\textbf{Área de aplicación:} Ingeniería Informática\\
\textbf{Director(es):} Dr. José Antonio Martínez García (\href{mailto:jamartine@ual.es}{jamartine@ual.es}) y Dr. Vicente González Ruiz (\href{mailto:vruiz@ual.es}{vruiz@ual.es}\texttt{)}

\begin{abstract}
	Se plantea el desarrollo de un sistema informático para la generación de imágenes estereoscópicas (por ejemplo, anaglifos) de la Luna a partir de tomas fotográficas capturadas al anochecer y al amanecer. Mediante el uso del paralaje generado por la propia rotación terrestre, el software a desarrollar aplicará algoritmos de tratamiento de imágenes para corregir la rotación de campo, escalar y alinear las imágenes, con el fin de obtener dicho esteograma (imagen en 3D). Diferentes formatos de visalización deberán generarse, incluyendo: esterogramas fotográficos (en color natural) y anaglifos (rojo y cian).
\end{abstract}

		
\section*{Planteamiento del Problema}
	
La percepción de profundidad en fotografía tridimensional requiere la captura de dos perspectivas de un mismo sujeto separadas por una distancia base (línea de base estereoscópica). En objetos astronómicos distantes como la Luna, resulta físicamente imposible obtener esta separación desplazando la cámara sobre la superficie terrestre en un mismo instante.
	
Sin embargo, es posible aprovechar el paralaje por rotación terrestre: durante un intervalo aproximado de 12 horas (desde el anochecer hasta el amanecer de una misma noche), la rotación de la Tierra desplaza al observador miles de kilómetros en el espacio, generando la línea de base necesaria para obtener un par estereoscópico natural.
	
El problema técnico principal reside en que las dos fotografías resultantes no son directamente fusionables. Debido al movimiento aparente de la Luna a lo largo de la bóveda celeste, las imágenes sufren severas discrepancias geométricas y radiométricas:
	
\begin{enumerate}
\item Rotación de campo: La orientación de la cara visible de la Luna cambia drásticamente en el encuadre.
\item Escala y perspectiva: Variaciones en la distancia focal, la refracción atmosférica cerca del horizonte y el tamaño aparente.
\item Cambios de iluminación y fase: Variaciones locales en las sombras de los cráteres y el contraste atmosférico durante la noche.
\end{enumerate}
	
Si estas diferencias no se corrigen algorítmicamente de forma previa, el sistema visual humano es incapaz de acoplar las imágenes, provocando fatiga ocular y la ruptura del efecto tridimensional.

\section*{Objetivos}

El objetivo principal es esarrollar un sistema informático de visión por computador y procesamiento digital de imágenes para la alineación, calibración y generación automática de autoestereogramas y estereogramas de la superficie lunar a partir de un par fotográfico capturado durante una misma noche (anochecer y amanecer). Más específicamente:

\begin{enumerate}
\item Diseñar e implementar un pipeline de procesamiento de imagen para el registro y alineación subpíxel del par de fotografías lunares.
\item Implementar algoritmos de corrección de rotación de campo, escalado y normalización de contraste/brillo.
\item Calcular el mapa de disparidad (mapa de profundidad) derivado del paralaje temporal entre ambas tomas.
\item Desarrollar el módulo de síntesis de autoestereogramas (formatos SIRDS/RDS y pares estereoscópicos/anaglifos) para la visualización 3D sin o con soporte óptico.
\item Evaluar la calidad de la reconstrucción tridimensional y la precisión del algoritmo de registro mediante métricas objetivas de disparidad.
\item Aplicar, opcionalmente, técnicas de superresolución.
\end{enumerate}
	
\section*{Resultados esperados}

Una aplicación que permita obtener un estereograma correcto de la Luna a partir de 2 imágenes de la misma tomadas tal y como se ha planteado.
	
\end{document}
